% !TEX root = ../main.tex
%-------------------------------------------------------------------------------
%	ABSTRACT PAGE
%-------------------------------------------------------------------------------

\addchaptertocentry{Хураангуй} % Хураангуйг гарчигт нэмэх

\begin{center}
{\scshape\Large \univname\par} % Их сургуулийн нэр
{\scshape\large \facname\par}\vspace{0.5cm} % Их сургуулийн нэр
{\huge\textbf{{Хураангуй}} \par}
\bigskip
{\Large{\ttitle} \par} % Тезисийн нэр
\bigskip

{\normalsize \shortname \par} % Зохиогчийн нэр
{\normalsize ariucn23@gmail.com \par}
\addressname

\end{center}

\textit{\textbf{Түлхүүр үгс: \keywordnames}}
\bigskip

Энэхүү дипломын ажлын хүрээнд рентген зургаас ясны хугарлыг ангилах, сегментлэх, хугарлын зэргийг тодорхойлох, эдгэрэлтийг таамаглах дөрвөн шаттай каскад гүн сургалтын pipeline боловсруулсан. Эхний хоёр туршилтад MURA өгөгдлийн санг зөвхөн суурь ангилал болон архитектурын туршилтад ашигласан. Харин эцсийн буюу гурав дахь туршилтын ангиллын модульд GRAZPEDWRI-DX, FracAtlas-COCO, YOLOBoneFracture болон BoneFractureCVProject өгөгдлийг нэгтгэн хэрэглэв. Сегментчлэлийн сургалтад FracAtlas-ийн маск болон YOLO хайрцаг тэмдэглэгээнээс SAM-аар үүсгэсэн псевдо маскуудыг ашигласан бол хугарлын зэргийн үнэлгээнд FracAtlas-ийн маскаас гаргасан морфологийн шинжүүдийг хэрэглэсэн. Эдгэрэлтийн модульд бодит follow-up шошго байхгүй тул синтетик клиникийн өгөгдлийг зөвхөн proof-of-concept түвшинд хэрэглэсэн.

Нэгдсэн олон даалгаварт загварт task interference-ийн системтэй асуудал илэрсний дараа бие даасан модульчлагдсан архитектурт шилжсэн. Эцсийн pipeline-ийн үр дүнгээр ангилагч (ConvNeXt-Base, ImageNet-22K + TTA) 3-fold cross-validation дээр AUC-ROC = 0.9887, F1 = 0.959; сегментчлэгч (Hybrid ResNet50 U-Net + TTA) тест дэх Dice = 0.604 гарсан нь анхны нэгдсэн загварын Dice = 0.0498-аас 12.1 дахин өндөр юм.

\textbf{Hold-out тестийн санд баталгаажуулалт}: Сургалтад огт ороогүй 574 зургийн нэгдсэн hold-out тестийн санд ангиллын Test AUC = 0.9703 (CV-аас generalization gap 0.018), F1 = 0.804; pipeline-conditional сегментчлэлийн Dice = 0.477 (TP=43). Үүнээс гадна, Stanford ML Group-ийн MURA-v1.1 өгөгдлийн санг гадаад (external) валидаци-д ашиглан 3{,}197 зураг дээр Test AUC = 0.8323 авч (XR\_FOREARM 0.93, XR\_HUMERUS 0.93, XR\_WRIST 0.87 хэсгүүдэд хамгийн өндөр), санал болгосон pipeline нь домэйнээс гадуурх өгөгдөл рүү ч ерөнхийлж чадахыг харуулсан.

Хугарлын зэргийн модуль нь морфологийн шинжүүдэд суурилан тест олонлог дээр Exact Accuracy = 1.000 үзүүлсэн боловч энэ үр дүн нь тест олонлогт зөвхөн Class~1--2 орсонтой холбоотой тул клиникийн баталгаажсан үр дүн бус, proof-of-concept хүрээний урьдчилсан үзүүлэлт гэж тайлбарлав. Загвар нь эмчийн хяналт дор клиникийн шийдвэрийг дэмжих (CDS) хэрэгсэл болох боломжтой боловч сегментчлэлийн тогтворжилт, хугарлын зэргийн клиник баталгаажуулалт болон эдгэрэлтийн бодит follow-up өгөгдөлд нэмэлт сайжруулалт шаардагдана.
